% ============================================================
% 00-sazetak-en.tex
% ============================================================

The class imbalance problem is one of the key challenges in machine
learning, particularly in classification tasks where the distribution of
examples across classes is not uniform. The SMOTE (Synthetic Minority
Over-sampling Technique) algorithm represents one of the most popular
approaches to address this problem by generating synthetic examples of
the minority class through linear interpolation.

This thesis analyzes the basic SMOTE algorithm and its variants:
Borderline-SMOTE, ADASYN, Safe-Level-SMOTE, K-Means SMOTE, SVM-SMOTE,
the combined approaches SMOTE-ENN and SMOTE-Tomek, and geometric
extensions G-SMOTE, Random SMOTE and Polynom-Fit SMOTE. Alternative
approaches such as cost-sensitive learning, ensemble methods, and neural
network applications are also considered.

The developed software solution implements the described algorithms in
Python and enables their experimental evaluation on a wide range of
datasets with varying characteristics. Statistical analysis of the
results was conducted to determine the behavior of individual algorithms
under different class ratios, dimensionalities, and noise levels.
Additionally, a web interface was developed for visualizing the synthetic
sample generation process.

\bigskip
\noindent\textbf{Keywords:} class imbalance, classification, oversampling,
SMOTE, synthetic samples
